\documentclass[aspectratio=169,10pt]{beamer}
\usetheme{Madrid}
\usecolortheme{default}
\usepackage[utf8]{inputenc}
\usepackage[T1]{fontenc}
\usepackage{lmodern}
\usepackage{amsmath,amssymb}
\usepackage{enumitem}
\usepackage{xcolor}

\setbeamertemplate{navigation symbols}{}
\setbeamertemplate{footline}[frame number]
\setbeamertemplate{itemize items}[circle]
\setbeamerfont{frametitle}{size=\large,series=\bfseries}

\title{CIA Part 1 -- Top 100 Hard Exam Questions}
\subtitle{Comprehensive Practice Deck: Governance, Risk, Controls, Ethics, IT \& QAIP}
\author{CIA Study Prep}
\date{\today}

\begin{document}
	
	\begin{frame}
		\titlepage
	\end{frame}
	
	\begin{frame}{Deck Overview}
		\begin{itemize}
			\item 100 curated high-level CIA Part 1 questions.
			\item Covers IIA Code of Ethics, Governance, Risk Management, Controls, Engagement Execution, IT \& Data Analytics, and QAIP.
			\item \textbf{Format:} Question slide followed by a detailed Solution/Explanation slide.
		\end{itemize}
	\end{frame}
	
	% ============ Q1 ============
	\begin{frame}[t]{Q1}
		\small
		\textbf{Question:} Which is the BEST description of "coordination" under Standard 2050?
		\begin{enumerate}[label=\Alph*.,itemsep=1pt,topsep=2pt]
			\item Sharing all workpapers with everyone.
			\item Coordinating activities with other assurance providers to ensure adequate coverage and minimize duplication.
			\item Avoiding contact with other assurance providers.
			\item Letting external audit direct internal audit.
		\end{enumerate}
	\end{frame}
	
	\begin{frame}[t]{Q1 -- Solution \& Explanation}
		\small
		\textbf{Correct Answer: B} \\[4pt]
		\textbf{Solution:} Standard 2050 requires coordination to ensure adequate coverage and minimize duplication. \\[4pt]
		\textbf{A wrong:} Unrestricted sharing may breach confidentiality. \\[2pt]
		\textbf{C wrong:} Avoidance causes duplication and gaps. \\[2pt]
		\textbf{D wrong:} External audit does not direct internal audit.
	\end{frame}
	
	% ============ Q2 ============
	\begin{frame}[t]{Q2}
		\small
		\textbf{Question:} Which is the MOST appropriate way for internal audit to rely on another assurance provider's work?
		\begin{enumerate}[label=\Alph*.,itemsep=1pt,topsep=2pt]
			\item Rely without evaluation.
			\item Evaluate the provider's competence, objectivity, and work quality before relying.
			\item Rely only if the provider is a personal contact.
			\item Rely only if it is inexpensive.
		\end{enumerate}
	\end{frame}
	
	\begin{frame}[t]{Q2 -- Solution \& Explanation}
		\small
		\textbf{Correct Answer: B} \\[4pt]
		\textbf{Solution:} Before relying on another provider, the auditor must evaluate competence, objectivity, and work quality. \\[4pt]
		\textbf{A wrong:} Blind reliance is improper and violates standards. \\[2pt]
		\textbf{C wrong:} Personal relationships are not an evaluation criterion. \\[2pt]
		\textbf{D wrong:} Cost is not a valid reliance criterion.
	\end{frame}
	
	% ============ Q3 ============
	\begin{frame}[t]{Q3}
		\small
		\textbf{Question:} Which is the BEST description of "reporting to senior management and the board" under Standard 2060?
		\begin{enumerate}[label=\Alph*.,itemsep=1pt,topsep=2pt]
			\item Reporting only at year-end.
			\item Reporting significant risk and control issues, including fraud risk and governance.
			\item Reporting only positive audit outcomes.
			\item Reporting findings directly to public media.
		\end{enumerate}
	\end{frame}
	
	\begin{frame}[t]{Q3 -- Solution \& Explanation}
		\small
		\textbf{Correct Answer: B} \\[4pt]
		\textbf{Solution:} Standard 2060 requires reporting significant risk and control issues, including fraud and governance. \\[4pt]
		\textbf{A wrong:} Reporting is ongoing as issues arise. \\[2pt]
		\textbf{C wrong:} Negative findings and adverse conditions must be reported. \\[2pt]
		\textbf{D wrong:} Media reporting breaches confidentiality.
	\end{frame}
	
	% ============ Q4 ============
	\begin{frame}[t]{Q4}
		\small
		\textbf{Question:} Which is the MOST appropriate way for internal audit to assess governance under Standard 2110?
		\begin{enumerate}[label=\Alph*.,itemsep=1pt,topsep=2pt]
			\item Evaluate the marketing plan.
			\item Evaluate ethics, board interaction, and governance processes.
			\item Prepare financial statements.
			\item Set strategic goals for management.
		\end{enumerate}
	\end{frame}
	
	\begin{frame}[t]{Q4 -- Solution \& Explanation}
		\small
		\textbf{Correct Answer: B} \\[4pt]
		\textbf{Solution:} Standard 2110 requires assessing governance, ethics, and board interaction processes. \\[4pt]
		\textbf{A wrong:} Operational marketing plans are not primary governance indicators. \\[2pt]
		\textbf{C wrong:} Financial statement preparation is management's responsibility. \\[2pt]
		\textbf{D wrong:} Setting strategy is management and the board's role.
	\end{frame}
	
	% ============ Q5 ============
	\begin{frame}[t]{Q5}
		\small
		\textbf{Question:} Which is the BEST description of "risk management" under Standard 2120?
		\begin{enumerate}[label=\Alph*.,itemsep=1pt,topsep=2pt]
			\item Preparing financial statements.
			\item Identifying, assessing, managing, and controlling risks.
			\item Setting enterprise strategy.
			\item Directing operational marketing.
		\end{enumerate}
	\end{frame}
	
	\begin{frame}[t]{Q5 -- Solution \& Explanation}
		\small
		\textbf{Correct Answer: B} \\[4pt]
		\textbf{Solution:} Risk management identifies, assesses, manages, and controls risk per Standard 2120. \\[4pt]
		\textbf{A wrong:} Financial reporting is an operational function. \\[2pt]
		\textbf{C wrong:} Setting strategy is management's task. \\[2pt]
		\textbf{D wrong:} Marketing is an operational domain.
	\end{frame}
	
	% ============ Q6 ============
	\begin{frame}[t]{Q6}
		\small
		\textbf{Question:} Which is the MOST appropriate way to assess the adequacy of an organization's control processes?
		\begin{enumerate}[label=\Alph*.,itemsep=1pt,topsep=2pt]
			\item Count total controls present.
			\item Evaluate whether controls mitigate risks within defined risk tolerance.
			\item Interview the CEO exclusively.
			\item Review the annual department budget.
		\end{enumerate}
	\end{frame}
	
	\begin{frame}[t]{Q6 -- Solution \& Explanation}
		\small
		\textbf{Correct Answer: B} \\[4pt]
		\textbf{Solution:} Controls are adequate if they effectively manage risks within established risk tolerances (Standard 2130). \\[4pt]
		\textbf{A wrong:} Quantity does not equal effectiveness. \\[2pt]
		\textbf{C wrong:} Single-source interviews are insufficient evidence. \\[2pt]
		\textbf{D wrong:} Budget amounts do not indicate control design or operation.
	\end{frame}
	
	% ============ Q7 ============
	\begin{frame}[t]{Q7}
		\small
		\textbf{Question:} Which is the MOST appropriate way to perform engagement-level risk assessment under Standard 2210.A3?
		\begin{enumerate}[label=\Alph*.,itemsep=1pt,topsep=2pt]
			\item Ignore operational risks.
			\item Assess risks relevant to the engagement's objectives and scope.
			\item Rely exclusively on last year's risk assessment.
			\item Use only management's risk register without testing.
		\end{enumerate}
	\end{frame}
	
	\begin{frame}[t]{Q7 -- Solution \& Explanation}
		\small
		\textbf{Correct Answer: B} \\[4pt]
		\textbf{Solution:} Engagement-level risk assessments must consider risks relevant to the specific engagement objectives (Standard 2210.A3). \\[4pt]
		\textbf{A wrong:} Ignoring risk violates IIA standards. \\[2pt]
		\textbf{C wrong:} Prior assessments must be re-validated. \\[2pt]
		\textbf{D wrong:} Management risk registers are inputs, not complete assessments.
	\end{frame}
	
	% ============ Q8 ============
	\begin{frame}[t]{Q8}
		\small
		\textbf{Question:} Which is the BEST description of "engagement resource allocation" under Standard 2230?
		\begin{enumerate}[label=\Alph*.,itemsep=1pt,topsep=2pt]
			\item Random allocation across engagements.
			\item Allocation based on the nature, complexity, and risk of the engagement.
			\item Equal allocation across all audit universe items.
			\item Allocation based on auditor tenure only.
		\end{enumerate}
	\end{frame}
	
	\begin{frame}[t]{Q8 -- Solution \& Explanation}
		\small
		\textbf{Correct Answer: B} \\[4pt]
		\textbf{Solution:} Resources must be allocated based on risk, complexity, and engagement nature (Standard 2230). \\[4pt]
		\textbf{A wrong:} Random allocation fails to address exposure. \\[2pt]
		\textbf{C wrong:} Equal allocation ignores risk variations. \\[2pt]
		\textbf{D wrong:} Seniority alone is not a risk-based metric.
	\end{frame}
	
	% ============ Q9 ============
	\begin{frame}[t]{Q9}
		\small
		\textbf{Question:} Which is the MOST appropriate way to evaluate information quality under Standard 2310?
		\begin{enumerate}[label=\Alph*.,itemsep=1pt,topsep=2pt]
			\item Assess sufficiency, reliability, relevance, and usefulness.
			\item Accept all management-provided information without testing.
			\item Ignore contradictory information.
			\item Utilize digital data exclusively.
		\end{enumerate}
	\end{frame}
	
	\begin{frame}[t]{Q9 -- Solution \& Explanation}
		\small
		\textbf{Correct Answer: A} \\[4pt]
		\textbf{Solution:} Standard 2310 requires auditors to collect sufficient, reliable, relevant, and useful information. \\[4pt]
		\textbf{B wrong:} Unverified acceptance violates due professional care. \\[2pt]
		\textbf{C wrong:} Contradictory evidence must be evaluated. \\[2pt]
		\textbf{D wrong:} Physical and oral evidence remain valid.
	\end{frame}
	
	% ============ Q10 ============
	\begin{frame}[t]{Q10}
		\small
		\textbf{Question:} Which is the BEST description of "analysis and evaluation" under Standard 2320?
		\begin{enumerate}[label=\Alph*.,itemsep=1pt,topsep=2pt]
			\item Gathering raw information.
			\item Analyzing and evaluating information to reach conclusions.
			\item Writing the final audit report.
			\item Setting organizational strategy.
		\end{enumerate}
	\end{frame}
	
	\begin{frame}[t]{Q10 -- Solution \& Explanation}
		\small
		\textbf{Correct Answer: B} \\[4pt]
		\textbf{Solution:} Standard 2320 requires auditors to base conclusions on appropriate analysis and evaluation. \\[4pt]
		\textbf{A wrong:} Information gathering falls under Standard 2310. \\[2pt]
		\textbf{C wrong:} Reporting is governed by Standard 2400. \\[2pt]
		\textbf{D wrong:} Setting strategy is a executive function.
	\end{frame}
	
	% ============ Q11 ============
	\begin{frame}[t]{Q11}
		\small
		\textbf{Question:} Which is the MOST appropriate way to evaluate "quality of communications" under Standard 2420?
		\begin{enumerate}[label=\Alph*.,itemsep=1pt,topsep=2pt]
			\item Ensure communications are detailed and highly complex.
			\item Ensure communications are accurate, objective, clear, concise, constructive, complete, and timely.
			\item Limit communications to verbal exit briefings.
			\item Use digital dashboard reporting exclusively.
		\end{enumerate}
	\end{frame}
	
	\begin{frame}[t]{Q11 -- Solution \& Explanation}
		\small
		\textbf{Correct Answer: B} \\[4pt]
		\textbf{Solution:} Standard 2420 explicitly lists seven quality criteria: accurate, objective, clear, concise, constructive, complete, and timely. \\[4pt]
		\textbf{A wrong:} Complexity hinders clarity. \\[2pt]
		\textbf{C wrong:} Formal written communications are frequently required. \\[2pt]
		\textbf{D wrong:} Communication formats should fit the audience.
	\end{frame}
	
	% ============ Q12 ============
	\begin{frame}[t]{Q12}
		\small
		\textbf{Question:} Which is the BEST description of "dissemination of results" under Standard 2440?
		\begin{enumerate}[label=\Alph*.,itemsep=1pt,topsep=2pt]
			\item Publicly sharing audit reports with external media.
			\item Distributing results to appropriate parties as defined by engagement objectives and charter.
			\item Posting internal reports on an open website.
			\item Restricting report dissemination solely to the auditee manager.
		\end{enumerate}
	\end{frame}
	
	\begin{frame}[t]{Q12 -- Solution \& Explanation}
		\small
		\textbf{Correct Answer: B} \\[4pt]
		\textbf{Solution:} Standard 2440 requires the CAE to communicate results to appropriate parties who can ensure proper evaluation. \\[4pt]
		\textbf{A wrong:} Public disclosure breaches confidentiality. \\[2pt]
		\textbf{C wrong:} Open posting compromises proprietary data. \\[2pt]
		\textbf{D wrong:} Governance bodies (board/management) must receive findings.
	\end{frame}
	
	% ============ Q13 ============
	\begin{frame}[t]{Q13}
		\small
		\textbf{Question:} Which is the MOST appropriate action if a communication contains confidential information that must be shared with the board?
		\begin{enumerate}[label=\Alph*.,itemsep=1pt,topsep=2pt]
			\item Distribute without restriction to all organization employees.
			\item Share with the board while explicitly limiting broader dissemination.
			\item Publish online with a disclaimer.
			\item Redact the information entirely before sending to the board.
		\end{enumerate}
	\end{frame}
	
	\begin{frame}[t]{Q13 -- Solution \& Explanation}
		\small
		\textbf{Correct Answer: B} \\[4pt]
		\textbf{Solution:} Confidential information may be shared with the board, but broader dissemination must be restricted (Standard 2440.A2). \\[4pt]
		\textbf{A wrong:} Unrestricted sharing breaches confidentiality policies. \\[2pt]
		\textbf{C wrong:} Publishing is improper. \\[2pt]
		\textbf{D wrong:} The board needs complete information to oversee risk.
	\end{frame}
	
	% ============ Q14 ============
	\begin{frame}[t]{Q14}
		\small
		\textbf{Question:} Which is the MOST appropriate way to perform ongoing internal QAIP assessments?
		\begin{enumerate}[label=\Alph*.,itemsep=1pt,topsep=2pt]
			\item Continuously evaluate engagement performance and operations.
			\item Conduct only one-time external reviews every five years.
			\item Review the annual audit budget exclusively.
			\item Issue random satisfaction surveys only.
		\end{enumerate}
	\end{frame}
	
	\begin{frame}[t]{Q14 -- Solution \& Explanation}
		\small
		\textbf{Correct Answer: A} \\[4pt]
		\textbf{Solution:} Ongoing internal assessments evaluate performance continuously across the audit activity (Standard 1311). \\[4pt]
		\textbf{A wrong:} Ongoing monitoring is continuous. \\[2pt]
		\textbf{B wrong:} External reviews are separate periodic requirements. \\[2pt]
		\textbf{C wrong:} Budget monitoring is administrative, not QAIP. \\[2pt]
		\textbf{D wrong:} Surveys represent only a single input.
	\end{frame}
	
	% ============ Q15 ============
	\begin{frame}[t]{Q15}
		\small
		\textbf{Question:} Which is the MOST appropriate action if an external quality assessment finds significant nonconformance?
		\begin{enumerate}[label=\Alph*.,itemsep=1pt,topsep=2pt]
			\item Ignore the findings.
			\item Develop an action plan and report the results to senior management and the board.
			\item Hide the results from stakeholders.
			\item Blame the external assessor for bias.
		\end{enumerate}
	\end{frame}
	
	\begin{frame}[t]{Q15 -- Solution \& Explanation}
		\small
		\textbf{Correct Answer: B} \\[4pt]
		\textbf{Solution:} Significant nonconformance requires corrective action plans and board reporting (Standard 1320). \\[4pt]
		\textbf{A wrong:} Ignoring findings violates IIA standards. \\[2pt]
		\textbf{C wrong:} Hiding results is professional misconduct. \\[2pt]
		\textbf{D wrong:} Blaming assessors is unprofessional.
	\end{frame}
	
	% ============ Q16 ============
	\begin{frame}[t]{Q16}
		\small
		\textbf{Question:} Which is the BEST description of "conformance" reporting under Standard 2430?
		\begin{enumerate}[label=\Alph*.,itemsep=1pt,topsep=2pt]
			\item Reporting total external audit fees.
			\item Reporting internal audit conformance with the Standards and Code of Ethics to the board.
			\item Reporting marketing metrics.
			\item Reporting departmental budget variances.
		\end{enumerate}
	\end{frame}
	
	\begin{frame}[t]{Q16 -- Solution \& Explanation}
		\small
		\textbf{Correct Answer: B} \\[4pt]
		\textbf{Solution:} The CAE reports conformance with the Standards and Code of Ethics to the board (Standard 2430). \\[4pt]
		\textbf{A wrong:} Audit fees are financial tracking. \\[2pt]
		\textbf{C wrong:} Marketing is unrelated. \\[2pt]
		\textbf{D wrong:} Budget variances are operational.
	\end{frame}
	
	% ============ Q17 ============
	\begin{frame}[t]{Q17}
		\small
		\textbf{Question:} Which is the MOST appropriate topic for the CAE to raise in a private session with the board?
		\begin{enumerate}[label=\Alph*.,itemsep=1pt,topsep=2pt]
			\item Strategic marketing plans.
			\item Sensitive issues management may be reluctant to address.
			\item Cafeteria operational changes.
			\item Individual staff auditor salary schedules.
		\end{enumerate}
	\end{frame}
	
	\begin{frame}[t]{Q17 -- Solution \& Explanation}
		\small
		\textbf{Correct Answer: B} \\[4pt]
		\textbf{Solution:} Private sessions give the CAE a secure forum to raise sensitive governance or management issues. \\[4pt]
		\textbf{A wrong:} Marketing is outside internal audit scope. \\[2pt]
		\textbf{C wrong:} Trivial administrative matters are inappropriate. \\[2pt]
		\textbf{D wrong:} Individual salaries are handled by administrative channels.
	\end{frame}
	
	% ============ Q18 ============
	\begin{frame}[t]{Q18}
		\small
		\textbf{Question:} Which is the MOST appropriate way to handle a scope restriction imposed by management?
		\begin{enumerate}[label=\Alph*.,itemsep=1pt,topsep=2pt]
			\item Accept the restriction without pushback.
			\item Report the restriction and its potential impact to the board per the charter.
			\item Ignore the restriction and proceed anyway.
			\item Quietly change engagement objectives to avoid conflict.
		\end{enumerate}
	\end{frame}
	
	\begin{frame}[t]{Q18 -- Solution \& Explanation}
		\small
		\textbf{Correct Answer: B} \\[4pt]
		\textbf{Solution:} Scope restrictions must be communicated to the board (Standard 1130/2600). \\[4pt]
		\textbf{A wrong:} Accepting restrictions compromises independence. \\[2pt]
		\textbf{C wrong:} Proceeding blindly creates legal and professional exposure. \\[2pt]
		\textbf{D wrong:} Objectives must not be altered silently.
	\end{frame}
	
	% ============ Q19 ============
	\begin{frame}[t]{Q19}
		\small
		\textbf{Question:} Which is the MOST likely red flag for fraudulent financial reporting?
		\begin{enumerate}[label=\Alph*.,itemsep=1pt,topsep=2pt]
			\item Revenue growth consistent with operating cash flows.
			\item Rapid revenue growth significantly exceeding cash flow generation.
			\item Strong, multi-layered internal controls.
			\item Timely regulatory filing disclosures.
		\end{enumerate}
	\end{frame}
	
	\begin{frame}[t]{Q19 -- Solution \& Explanation}
		\small
		\textbf{Correct Answer: B} \\[4pt]
		\textbf{Solution:} Revenue growing much faster than cash flow indicates potential premature or fictitious revenue recognition. \\[4pt]
		\textbf{A wrong:} Consistency is a normal business indicator. \\[2pt]
		\textbf{C wrong:} Strong controls lower fraud risk. \\[2pt]
		\textbf{D wrong:} Timely filing is a positive governance sign.
	\end{frame}
	
	% ============ Q20 ============
	\begin{frame}[t]{Q20}
		\small
		\textbf{Question:} Which is the MOST appropriate action if internal audit is asked to participate in a suspected fraud investigation?
		\begin{enumerate}[label=\Alph*.,itemsep=1pt,topsep=2pt]
			\item Interrogate suspects independently.
			\item Support the investigation by gathering evidence while legal/management lead.
			\item Make the final legal determination of guilt.
			\item Impose disciplinary employment actions.
		\end{enumerate}
	\end{frame}
	
	\begin{frame}[t]{Q20 -- Solution \& Explanation}
		\small
		\textbf{Correct Answer: B} \\[4pt]
		\textbf{Solution:} Internal audit supports fraud investigations through evidence gathering, while legal and management lead the process. \\[4pt]
		\textbf{A wrong:} Interrogation belongs to legal, security, or HR. \\[2pt]
		\textbf{C wrong:} Determining guilt is a judicial role. \\[2pt]
		\textbf{D wrong:} Discipline is management's responsibility.
	\end{frame}
	
	% ============ Q21 ============
	\begin{frame}[t]{Q21}
		\small
		\textbf{Question:} Which is the BEST description of "fraud risk factors"?
		\begin{enumerate}[label=\Alph*.,itemsep=1pt,topsep=2pt]
			\item Conditions or events indicating incentives, pressures, or opportunities to commit fraud.
			\item The calculated internal audit fee.
			\item The departmental operating budget.
			\item Corporate marketing strategy documents.
		\end{enumerate}
	\end{frame}
	
	\begin{frame}[t]{Q21 -- Solution \& Explanation}
		\small
		\textbf{Correct Answer: A} \\[4pt]
		\textbf{Solution:} Fraud risk factors are conditions reflecting incentives, opportunities, or rationalizations. \\[4pt]
		\textbf{B wrong:} Fees are financial transactions. \\[2pt]
		\textbf{C wrong:} Budgets are financial plans. \\[2pt]
		\textbf{D wrong:} Marketing is an operational area.
	\end{frame}
	
	% ============ Q22 ============
	\begin{frame}[t]{Q22}
		\small
		\textbf{Question:} Which is the BEST description of a "process owner" under the IIA Glossary?
		\begin{enumerate}[label=\Alph*.,itemsep=1pt,topsep=2pt]
			\item The individual or group accountable for the operational process being evaluated.
			\item The senior internal auditor performing fieldwork.
			\item The audit committee chair.
			\item The external statutory auditor.
		\end{enumerate}
	\end{frame}
	
	\begin{frame}[t]{Q22 -- Solution \& Explanation}
		\small
		\textbf{Correct Answer: A} \\[4pt]
		\textbf{Solution:} A process owner is responsible for managing and maintaining the activity being reviewed. \\[4pt]
		\textbf{B wrong:} The internal auditor performs evaluation, not ownership. \\[2pt]
		\textbf{C wrong:} The audit committee provides oversight. \\[2pt]
		\textbf{D wrong:} External auditors provide independent financial assurance.
	\end{frame}
	
	% ============ Q23 ============
	\begin{frame}[t]{Q23}
		\small
		\textbf{Question:} Which is the BEST description of an "auditee" under the IIA Glossary?
		\begin{enumerate}[label=\Alph*.,itemsep=1pt,topsep=2pt]
			\item The person, group, or function subject to an internal audit engagement.
			\item The lead internal auditor.
			\item The board of directors.
			\item The external audit engagement partner.
		\end{enumerate}
	\end{frame}
	
	\begin{frame}[t]{Q23 -- Solution \& Explanation}
		\small
		\textbf{Correct Answer: A} \\[4pt]
		\textbf{Solution:} The auditee is the specific entity or function being audited. \\[4pt]
		\textbf{A wrong:} The auditee is the subject of review. \\[2pt]
		\textbf{B wrong:} The auditor conducts the review. \\[2pt]
		\textbf{C wrong:} The board oversees the function. \\[2pt]
		\textbf{D wrong:} External partners lead independent audits.
	\end{frame}
	
	% ============ Q24 ============
	\begin{frame}[t]{Q24}
		\small
		\textbf{Question:} Which is the BEST description of a "stakeholder" under the IIA Glossary?
		\begin{enumerate}[label=\Alph*.,itemsep=1pt,topsep=2pt]
			\item Anyone with a vested interest or concern in the organization's success.
			\item The lead staff internal auditor.
			\item The departmental expense budget.
			\item The finalized internal audit plan.
		\end{enumerate}
	\end{frame}
	
	\begin{frame}[t]{Q24 -- Solution \& Explanation}
		\small
		\textbf{Correct Answer: A} \\[4pt]
		\textbf{Solution:} Stakeholders include all parties with an interest in the organization's objectives (board, management, regulators, public). \\[4pt]
		\textbf{B wrong:} Auditors are participants, not general stakeholders. \\[2pt]
		\textbf{C wrong:} Budgets are financial management tools. \\[2pt]
		\textbf{D wrong:} Plans are operational documents.
	\end{frame}
	
	% ============ Q25 ============
	\begin{frame}[t]{Q25}
		\small
		\textbf{Question:} Which is the BEST description of a "process" under the IIA Glossary?
		\begin{enumerate}[label=\Alph*.,itemsep=1pt,topsep=2pt]
			\item A set of interrelated or interacting activities that transforms inputs into outputs.
			\item The final printed audit report.
			\item The annual financial budget.
			\item The corporate marketing outline.
		\end{enumerate}
	\end{frame}
	
	\begin{frame}[t]{Q25 -- Solution \& Explanation}
		\small
		\textbf{Correct Answer: A} \\[4pt]
		\textbf{Solution:} A process is an interconnected set of activities designed to achieve specific objectives. \\[4pt]
		\textbf{B wrong:} The report is an output deliverable. \\[2pt]
		\textbf{C wrong:} Budgets are financial controls. \\[2pt]
		\textbf{D wrong:} Marketing plans are operational strategies.
	\end{frame}
	
	% ============ Q26 ============
	\begin{frame}[t]{Q26}
		\small
		\textbf{Question:} Which is the MOST appropriate way to assess the adequacy of internal audit resource levels?
		\begin{enumerate}[label=\Alph*.,itemsep=1pt,topsep=2pt]
			\item Count total audit headcount relative to industry averages.
			\item Assess whether available resources allow successful execution of the approved internal audit plan.
			\item Review the total dollar allocation of the travel budget.
			\item Conduct brief interviews with corporate human resources.
		\end{enumerate}
	\end{frame}
	
	\begin{frame}[t]{Q26 -- Solution \& Explanation}
		\small
		\textbf{Correct Answer: B} \\[4pt]
		\textbf{Solution:} Resource adequacy is measured by the ability to fulfill the internal audit mandate and plan (Standard 2030). \\[4pt]
		\textbf{A wrong:} Raw headcount ignores skill sets and risk coverage. \\[2pt]
		\textbf{C wrong:} Travel budgets do not measure total resource capacity. \\[2pt]
		\textbf{D wrong:} HR interviews do not measure audit coverage capability.
	\end{frame}
	
	% ============ Q27 ============
	\begin{frame}[t]{Q27}
		\small
		\textbf{Question:} Which is the BEST description of audit committee oversight responsibilities?
		\begin{enumerate}[label=\Alph*.,itemsep=1pt,topsep=2pt]
			\item Performing direct fieldwork testing.
			\item Overseeing internal audit independence, performance quality, and reporting lines.
			\item Writing detailed engagement workpapers.
			\item Negotiating commercial vendor contracts.
		\end{enumerate}
	\end{frame}
	
	\begin{frame}[t]{Q27 -- Solution \& Explanation}
		\small
		\textbf{Correct Answer: B} \\[4pt]
		\textbf{Solution:} The audit committee provides governance oversight over internal audit's independence, performance, and results. \\[4pt]
		\textbf{A wrong:} Internal audit executes fieldwork, not the committee. \\[2pt]
		\textbf{C wrong:} Workpapers are compiled by internal auditors. \\[2pt]
		\textbf{D wrong:} Vendor contracting is an executive management duty.
	\end{frame}
	
	% ============ Q28 ============
	\begin{frame}[t]{Q28}
		\small
		\textbf{Question:} Which test is MOST appropriate for verifying the completeness of recorded accounts payable?
		\begin{enumerate}[label=\Alph*.,itemsep=1pt,topsep=2pt]
			\item Vouching entries from the general ledger down to vendor invoices.
			\item Tracing from unrecorded receiving reports and vendor invoices up to the general ledger.
			\item Confirming high-balance payables with major customers.
			\item Reviewing the corporate marketing budget.
		\end{enumerate}
	\end{frame}
	
	\begin{frame}[t]{Q28 -- Solution \& Explanation}
		\small
		\textbf{Correct Answer: B} \\[4pt]
		\textbf{Solution:} Tracing source documents (receiving reports/invoices) to the ledger tests completeness. \\[4pt]
		\textbf{A wrong:} Ledger-to-invoice vouching tests existence. \\[2pt]
		\textbf{C wrong:} Customers confirm receivables, not payables. \\[2pt]
		\textbf{D wrong:} Budgets are irrelevant to payable completeness.
	\end{frame}
	
	% ============ Q29 ============
	\begin{frame}[t]{Q29}
		\small
		\textbf{Question:} Which is the MOST likely effect of an effective internal audit activity on corporate governance?
		\begin{enumerate}[label=\Alph*.,itemsep=1pt,topsep=2pt]
			\item It weakens governance frameworks by duplicating duties.
			\item It strengthens governance by providing independent objective assurance and insight.
			\item It has zero measurable impact on organizational tone.
			\item It replaces the board of directors entirely.
		\end{enumerate}
	\end{frame}
	
	\begin{frame}[t]{Q29 -- Solution \& Explanation}
		\small
		\textbf{Correct Answer: B} \\[4pt]
		\textbf{Solution:} Internal audit strengthens governance via objective assurance and advice. \\[4pt]
		\textbf{A wrong:} Effective coordination prevents duplication. \\[2pt]
		\textbf{C wrong:} It has a substantial positive impact. \\[2pt]
		\textbf{D wrong:} Internal audit supports, but never replaces, the board.
	\end{frame}
	
	% ============ Q30 ============
	\begin{frame}[t]{Q30}
		\small
		\textbf{Question:} Which control procedure is MOST effective for verifying the accuracy of a vendor invoice prior to disbursement?
		\begin{enumerate}[label=\Alph*.,itemsep=1pt,topsep=2pt]
			\item Obtaining a single signature from the accounts payable clerk.
			\item Performing a three-way match (Purchase Order, Receiving Report, and Vendor Invoice).
			\item Making an informal telephone call to the vendor sales rep.
			\item Reviewing the departmental annual financial budget.
		\end{enumerate}
	\end{frame}
	
	\begin{frame}[t]{Q30 -- Solution \& Explanation}
		\small
		\textbf{Correct Answer: B} \\[4pt]
		\textbf{Solution:} The three-way match verifies quantities, prices, and receipt of goods before payment. \\[4pt]
		\textbf{A wrong:} A single signature lacks verification controls. \\[2pt]
		\textbf{C wrong:} Informal calls lack audit trails. \\[2pt]
		\textbf{D wrong:} Budget review does not verify transaction accuracy.
	\end{frame}
	
	% ============ Q31 ============
	\begin{frame}[t]{Q31}
		\small
		\textbf{Question:} Which is the BEST description of "materiality" in internal auditing?
		\begin{enumerate}[label=\Alph*.,itemsep=1pt,topsep=2pt]
			\item A fixed 5 percent threshold applied universally.
			\item The magnitude of an omission or misstatement that could influence the decisions of a reasonable user.
			\item The total annual audit fee.
			\item The total operating budget allocation.
		\end{enumerate}
	\end{frame}
	
	\begin{frame}[t]{Q31 -- Solution \& Explanation}
		\small
		\textbf{Correct Answer: B} \\[4pt]
		\textbf{Solution:} Materiality evaluates whether an omission or misstatement would influence a reasonable user's judgment. \\[4pt]
		\textbf{A wrong:} Materiality involves qualitative and quantitative professional judgment, not a rigid fixed percent. \\[2pt]
		\textbf{C wrong:} Fees are administrative. \\[2pt]
		\textbf{D wrong:} Budgets are resource plans.
	\end{frame}
	
	% ============ Q32 ============
	\begin{frame}[t]{Q32}
		\small
		\textbf{Question:} Which is the MOST appropriate action if an auditor discovers that auditee personnel intentionally destroyed records?
		\begin{enumerate}[label=\Alph*.,itemsep=1pt,topsep=2pt]
			\item Attempt to reconstruct the records independently without documenting the incident.
			\item Document the spoliation of evidence thoroughly and notify legal counsel and the board.
			\item Ignore the missing documents since testing can proceed around them.
			\item Confront the auditee aggressively in public to force a confession.
		\end{enumerate}
	\end{frame}
	
	\begin{frame}[t]{Q32 -- Solution \& Explanation}
		\small
		\textbf{Correct Answer: B} \\[4pt]
		\textbf{Solution:} Record destruction (spoliation) is a severe legal and control violation; document and escalate. \\[4pt]
		\textbf{A wrong:} Independent reconstruction hides illegal evidence tampering. \\[2pt]
		\textbf{C wrong:} Ignoring missing files violates due care. \\[2pt]
		\textbf{D wrong:} Aggressive confrontation increases risk and breaches professionalism.
	\end{frame}
	
	% ============ Q33 ============
	\begin{frame}[t]{Q33}
		\small
		\textbf{Question:} Which is the BEST description of a "control deficiency"?
		\begin{enumerate}[label=\Alph*.,itemsep=1pt,topsep=2pt]
			\item A control that is overly expensive to operate.
			\item A condition where a control is missing or does not operate as designed to prevent/detect misstatements timely.
			\item A control that lacks physical documentation.
			\item An automated IT application control.
		\end{enumerate}
	\end{frame}
	
	\begin{frame}[t]{Q33 -- Solution \& Explanation}
		\small
		\textbf{Correct Answer: B} \\[4pt]
		\textbf{Solution:} A control deficiency occurs when a control fails to operate effectively or is missing. \\[4pt]
		\textbf{A wrong:} High cost is an efficiency issue, not a design deficiency. \\[2pt]
		\textbf{C wrong:} Documentation gaps are only one form of deficiency. \\[2pt]
		\textbf{D wrong:} Automated controls can be effective or deficient.
	\end{frame}
	
	% ============ Q34 ============
	\begin{frame}[t]{Q34}
		\small
		\textbf{Question:} Which is the MOST appropriate method for testing logical access controls within an ERP system?
		\begin{enumerate}[label=\Alph*.,itemsep=1pt,topsep=2pt]
			\item Review the annual IT departmental expense budget.
			\item Compare actual user system access rights against the approved user access matrix.
			\item Interview the Chief Information Officer about security policies.
			\item Review external firewall hardware configuration guides.
		\end{enumerate}
	\end{frame}
	
	\begin{frame}[t]{Q34 -- Solution \& Explanation}
		\small
		\textbf{Correct Answer: B} \\[4pt]
		\textbf{Solution:} Comparing user access lists against an approved segregation matrix tests logical access control. \\[4pt]
		\textbf{A wrong:} Budgets do not test system security. \\[2pt]
		\textbf{C wrong:} Inquiry alone provides weak evidence. \\[2pt]
		\textbf{D wrong:} Firewalls are perimeter controls, not ERP application user access.
	\end{frame}
	
	% ============ Q35 ============
	\begin{frame}[t]{Q35}
		\small
		\textbf{Question:} Which is the BEST description of a "material weakness"?
		\begin{enumerate}[label=\Alph*.,itemsep=1pt,topsep=2pt]
			\item A minor clerical error in accounting records.
			\item A deficiency (or combination of deficiencies) resulting in a reasonable possibility of a material misstatement.
			\item A missing organizational policy manual.
			\item An employee training deficit.
		\end{enumerate}
	\end{frame}
	
	\begin{frame}[t]{Q35 -- Solution \& Explanation}
		\small
		\textbf{Correct Answer: B} \\[4pt]
		\textbf{Solution:} A material weakness carries a reasonable possibility of resulting in a material misstatement. \\[4pt]
		\textbf{A wrong:} Minor errors are not material. \\[2pt]
		\textbf{C wrong:} Policy gaps may or may not be material. \\[2pt]
		\textbf{D wrong:} Training needs do not automatically equal material weaknesses.
	\end{frame}
	
	% ============ Q36 ============
	\begin{frame}[t]{Q36}
		\small
		\textbf{Question:} Which is the MOST appropriate action if an auditor discovers clear evidence of management override of controls?
		\begin{enumerate}[label=\Alph*.,itemsep=1pt,topsep=2pt]
			\item Ignore the finding if the financial impact is minimal.
			\item Expand substantive testing procedures and report the override to governance oversight.
			\item Reduce substantive audit testing due to management cooperation.
			\item Rely entirely on management representations.
		\end{enumerate}
	\end{frame}
	
	\begin{frame}[t]{Q36 -- Solution \& Explanation}
		\small
		\textbf{Correct Answer: B} \\[4pt]
		\textbf{Solution:} Management override represents high risk, requiring expanded testing and escalation. \\[4pt]
		\textbf{A wrong:} Override is inherently significant and cannot be ignored. \\[2pt]
		\textbf{C wrong:} Testing must be increased, not reduced. \\[2pt]
		\textbf{D wrong:} Management representations cannot replace objective testing when override occurs.
	\end{frame}
	
	% ============ Q37 ============
	\begin{frame}[t]{Q37}
		\small
		\textbf{Question:} Which is the BEST description of "reasonable assurance" in internal auditing?
		\begin{enumerate}[label=\Alph*.,itemsep=1pt,topsep=2pt]
			\item Absolute, 100 percent certainty that no errors exist.
			\item A high, but not absolute, level of assurance constrained by inherent audit limitations.
			\item Zero assurance provided to management.
			\item Management's personal performance guarantee.
		\end{enumerate}
	\end{frame}
	
	\begin{frame}[t]{Q37 -- Solution \& Explanation}
		\small
		\textbf{Correct Answer: B} \\[4pt]
		\textbf{Solution:} Reasonable assurance is high-level assurance limited by evidence constraints and cost-benefit balances. \\[4pt]
		\textbf{A wrong:} Absolute certainty is impossible in auditing. \\[2pt]
		\textbf{C wrong:} Internal audit provides positive assurance. \\[2pt]
		\textbf{D wrong:} Management cannot guarantee controls.
	\end{frame}
	
	% ============ Q38 ============
	\begin{frame}[t]{Q38}
		\small
		\textbf{Question:} Which is the MOST appropriate action when an auditor identifies indicators of possible fraud?
		\begin{enumerate}[label=\Alph*.,itemsep=1pt,topsep=2pt]
			\item Conduct full interrogations of suspected employees alone.
			\item Report indicators to appropriate management or the board and recommend a formal investigation.
			\item Defer reporting until the final scheduled year-end audit report.
			\item Conceal the indicators to protect employee morale.
		\end{enumerate}
	\end{frame}
	
	\begin{frame}[t]{Q38 -- Solution \& Explanation}
		\small
		\textbf{Correct Answer: B} \\[4pt]
		\textbf{Solution:} Auditors report fraud indicators and recommend investigations; they do not lead criminal interrogations. \\[4pt]
		\textbf{A wrong:} Solo interrogations compromise legal processes. \\[2pt]
		\textbf{C wrong:} Fraud indicators require prompt notification. \\[2pt]
		\textbf{D wrong:} Concealing fraud violates professional ethics.
	\end{frame}
	
	% ============ Q39 ============
	\begin{frame}[t]{Q39}
		\small
		\textbf{Question:} Which is the BEST description of the "control environment"?
		\begin{enumerate}[label=\Alph*.,itemsep=1pt,topsep=2pt]
			\item Network firewall hardware settings.
			\item The overall attitude, awareness, and actions of management and the board regarding ethics and internal control.
			\item The formal annual audit plan document.
			\item The external audit firm's opinion letter.
		\end{enumerate}
	\end{frame}
	
	\begin{frame}[t]{Q39 -- Solution \& Explanation}
		\small
		\textbf{Correct Answer: B} \\[4pt]
		\textbf{Solution:} The control environment represents the tone at the top regarding ethics and internal control. \\[4pt]
		\textbf{A wrong:} Firewalls are specific control activities. \\[2pt]
		\textbf{C wrong:} The plan is an operational schedule. \\[2pt]
		\textbf{D wrong:} External opinions are outside internal governance.
	\end{frame}
	
	% ============ Q40 ============
	\begin{frame}[t]{Q40}
		\small
		\textbf{Question:} Which is the MOST appropriate response when an auditor is assigned to evaluate an operational activity they managed 10 months prior?
		\begin{enumerate}[label=\Alph*.,itemsep=1pt,topsep=2pt]
			\item Proceed with the audit normally since sufficient time has elapsed.
			\item Disclose the impairment of objectivity to management and request reassignment.
			\item Proceed with the audit but use a smaller sample size.
			\item Ignore the prior management role entirely.
		\end{enumerate}
	\end{frame}
	
	\begin{frame}[t]{Q40 -- Solution \& Explanation}
		\small
		\textbf{Correct Answer: B} \\[4pt]
		\textbf{Solution:} Prior management responsibilities impair objectivity; auditors must disclose this and seek reassignment (Standard 1130). \\[4pt]
		\textbf{A wrong:} Time passage does not automatically cure impairment within one year. \\[2pt]
		\textbf{C wrong:} Sample size adjustments do not restore objectivity. \\[2pt]
		\textbf{D wrong:} Prior roles cannot be concealed.
	\end{frame}
	
	% ============ Q41 ============
	\begin{frame}[t]{Q41}
		\small
		\textbf{Question:} Which audit procedure provides the BEST evidence for verifying the physical existence of inventory?
		\begin{enumerate}[label=\Alph*.,itemsep=1pt,topsep=2pt]
			\item Reviewing perpetual inventory ledger balances.
			\item Direct physical observation and inspection during inventory counts.
			\item Conducting an interview with the warehouse supervisor.
			\item Inspecting the annual inventory purchase budget.
		\end{enumerate}
	\end{frame}
	
	\begin{frame}[t]{Q41 -- Solution \& Explanation}
		\small
		\textbf{Correct Answer: B} \\[4pt]
		\textbf{Solution:} Physical observation provides primary physical evidence of existence. \\[4pt]
		\textbf{A wrong:} Perpetual records are internal documentation. \\[2pt]
		\textbf{C wrong:} Inquiry alone provides weak evidence. \\[2pt]
		\textbf{D wrong:} Budgets do not confirm physical existence.
	\end{frame}
	
	% ============ Q42 ============
	\begin{frame}[t]{Q42}
		\small
		\textbf{Question:} Which is the BEST description of "residual risk"?
		\begin{enumerate}[label=\Alph*.,itemsep=1pt,topsep=2pt]
			\item The risk of a misstatement remaining undetected by the auditor.
			\item The risk remaining after management takes action to alter risk likelihood and impact.
			\item The risk of fraud occurring before any controls are applied.
			\item The risk associated with engagement staffing shortages.
		\end{enumerate}
	\end{frame}
	
	\begin{frame}[t]{Q42 -- Solution \& Explanation}
		\small
		\textbf{Correct Answer: B} \\[4pt]
		\textbf{Solution:} Residual risk is the risk remaining after management implements risk responses and controls. \\[4pt]
		\textbf{A wrong:} Undetected misstatement risk is detection risk. \\[2pt]
		\textbf{C wrong:} Risk before controls is inherent risk. \\[2pt]
		\textbf{D wrong:} Staffing issues relate to resource risk.
	\end{frame}
	
	% ============ Q43 ============
	\begin{frame}[t]{Q43}
		\small
		\textbf{Question:} Which is the MOST appropriate way to coordinate internal audit activities with external auditors?
		\begin{enumerate}[label=\Alph*.,itemsep=1pt,topsep=2pt]
			\item Share all internal audit files unrestricted without reviewing confidentiality protocols.
			\item Coordinate timing, scope, and access to workpapers while protecting independence.
			\item Allow external auditors to dictate internal audit's annual plan and budget.
			\item Avoid all direct contact to maintain absolute separation.
		\end{enumerate}
	\end{frame}
	
	\begin{frame}[t]{Q43 -- Solution \& Explanation}
		\small
		\textbf{Correct Answer: B} \\[4pt]
		\textbf{Solution:} Coordination should align timing and scope while maintaining internal audit independence. \\[4pt]
		\textbf{A wrong:} Unrestricted sharing may violate confidentiality rules. \\[2pt]
		\textbf{C wrong:} External audit must not direct internal audit. \\[2pt]
		\textbf{D wrong:} Complete avoidance results in inefficient overlap.
	\end{frame}
	
	% ============ Q44 ============
	\begin{frame}[t]{Q44}
		\small
		\textbf{Question:} Which elements comprise the classic "fraud triangle"?
		\begin{enumerate}[label=\Alph*.,itemsep=1pt,topsep=2pt]
			\item Inherent risk, control risk, and detection risk.
			\item Pressure/incentive, opportunity, and rationalization.
			\item Plan, do, check, and act.
			\item Cash, inventory, and accounts receivable.
		\end{enumerate}
	\end{frame}
	
	\begin{frame}[t]{Q44 -- Solution \& Explanation}
		\small
		\textbf{Correct Answer: B} \\[4pt]
		\textbf{Solution:} Cressey's fraud triangle consists of pressure/incentive, opportunity, and rationalization. \\[4pt]
		\textbf{A wrong:} Those form the audit risk model. \\[2pt]
		\textbf{C wrong:} That is the continuous improvement (PDCA) cycle. \\[2pt]
		\textbf{D wrong:} Those are financial balance sheet accounts.
	\end{frame}
	
	% ============ Q45 ============
	\begin{frame}[t]{Q45}
		\small
		\textbf{Question:} Which audit test provides the BEST evidence regarding the operating effectiveness of segregation of duties?
		\begin{enumerate}[label=\Alph*.,itemsep=1pt,topsep=2pt]
			\item Observing a single transaction execution once during fieldwork.
			\item Examining system user access rights, profiles, and role configurations.
			\item Conducting a brief interview with the department manager.
			\item Reviewing the organization chart graphic.
		\end{enumerate}
	\end{frame}
	
	\begin{frame}[t]{Q45 -- Solution \& Explanation}
		\small
		\textbf{Correct Answer: B} \\[4pt]
		\textbf{Solution:} Examining actual system access rights and roles provides strong evidence of segregation. \\[4pt]
		\textbf{A wrong:} Single observations provide limited coverage. \\[2pt]
		\textbf{C wrong:} Inquiry alone is insufficient. \\[2pt]
		\textbf{D wrong:} Org charts do not show system authorization rights.
	\end{frame}
	
	% ============ Q46 ============
	\begin{frame}[t]{Q46}
		\small
		\textbf{Question:} Which is the BEST description of "inherent risk"?
		\begin{enumerate}[label=\Alph*.,itemsep=1pt,topsep=2pt]
			\item The risk remaining after management implements controls.
			\item The susceptibility of an information item or process to error or fraud assuming no related controls exist.
			\item The risk that the auditor fails to uncover a material error.
			\item The risk of corporate legal noncompliance.
		\end{enumerate}
	\end{frame}
	
	\begin{frame}[t]{Q46 -- Solution \& Explanation}
		\small
		\textbf{Correct Answer: B} \\[4pt]
		\textbf{Solution:} Inherent risk is the exposure existing prior to considering internal controls. \\[4pt]
		\textbf{A wrong:} Post-control risk is residual risk. \\[2pt]
		\textbf{C wrong:} Auditor failure risk is detection risk. \\[2pt]
		\textbf{D wrong:} Compliance risk is a specific category of risk.
	\end{frame}
	
	% ============ Q47 ============
	\begin{frame}[t]{Q47}
		\small
		\textbf{Question:} Which is the BEST description of "detection risk"?
		\begin{enumerate}[label=\Alph*.,itemsep=1pt,topsep=2pt]
			\item The risk that internal audit procedures fail to detect a material misstatement.
			\item The risk that internal controls fail to prevent a misstatement.
			\item The risk inherent in the operational environment.
			\item The risk of executive management fraud override.
		\end{enumerate}
	\end{frame}
	
	\begin{frame}[t]{Q47 -- Solution \& Explanation}
		\small
		\textbf{Correct Answer: A} \\[4pt]
		\textbf{Solution:} Detection risk is the risk that audit procedures fail to identify existing errors. \\[4pt]
		\textbf{B wrong:} Control failure risk is control risk. \\[2pt]
		\textbf{C wrong:} Environmental risk is inherent risk. \\[2pt]
		\textbf{D wrong:} Override relates to control environment and fraud risk.
	\end{frame}
	
	% ============ Q48 ============
	\begin{frame}[t]{Q48}
		\small
		\textbf{Question:} Which procedure is MOST appropriate for testing a routine bank reconciliation control?
		\begin{enumerate}[label=\Alph*.,itemsep=1pt,topsep=2pt]
			\item Reading the written accounting policy manual.
			\item Re-performing the reconciliation calculation and investigating reconciling items.
			\item Interviewing the accounts payable supervisor.
			\item Reviewing the annual cash flow forecast.
		\end{enumerate}
	\end{frame}
	
	\begin{frame}[t]{Q48 -- Solution \& Explanation}
		\small
		\textbf{Correct Answer: B} \\[4pt]
		\textbf{Solution:} Re-performing reconciliations tests whether the control operates as intended. \\[4pt]
		\textbf{A wrong:} Policy review tests design, not operation. \\[2pt]
		\textbf{C wrong:} Inquiry alone provides weak evidence. \\[2pt]
		\textbf{D wrong:} Cash flow forecasts are forward-looking budgets.
	\end{frame}
	
	% ============ Q49 ============
	\begin{frame}[t]{Q49}
		\small
		\textbf{Question:} Which is the MOST appropriate action when an auditor identifies an unmitigated conflict of interest?
		\begin{enumerate}[label=\Alph*.,itemsep=1pt,topsep=2pt]
			\item Conceal the conflict from audit leadership.
			\item Document the conflict and escalate through proper organizational channels.
			\item Confront the individual involved in a hostile manner.
			\item Leak details of the conflict to external news media.
		\end{enumerate}
	\end{frame}
	
	\begin{frame}[t]{Q49 -- Solution \& Explanation}
		\small
		\textbf{Correct Answer: B} \\[4pt]
		\textbf{Solution:} Conflicts of interest must be documented and reported through formal channels. \\[4pt]
		\textbf{A wrong:} Concealment violates the Code of Ethics. \\[2pt]
		\textbf{C wrong:} Hostile confrontation increases organizational risk. \\[2pt]
		\textbf{D wrong:} Leaking information breaches confidentiality.
	\end{frame}
	
	% ============ Q50 ============
	\begin{frame}[t]{Q50}
		\small
		\textbf{Question:} Which practices represent foundational elements of basic "cyber hygiene"?
		\begin{enumerate}[label=\Alph*.,itemsep=1pt,topsep=2pt]
			\item Corporate marketing strategy alignment.
			\item Regular software patching, secure backups, and strict access controls.
			\item Minimizing physical office desk space.
			\item Reducing internal audit reporting frequency.
		\end{enumerate}
	\end{frame}
	
	\begin{frame}[t]{Q50 -- Solution \& Explanation}
		\small
		\textbf{Correct Answer: B} \\[4pt]
		\textbf{Solution:} Cyber hygiene covers essential IT security practices such as patching, backups, and access management. \\[4pt]
		\textbf{A wrong:} Marketing is unrelated to IT security. \\[2pt]
		\textbf{C wrong:} Office layout does not constitute cyber hygiene. \\[2pt]
		\textbf{D wrong:} Reporting frequency does not affect security.
	\end{frame}
	
	% ============ Q51 ============
	\begin{frame}[t]{Q51}
		\small
		\textbf{Question:} Which is the FIRST step in planning an internal audit engagement under Standard 2200?
		\begin{enumerate}[label=\Alph*.,itemsep=1pt,topsep=2pt]
			\item Performing substantive transaction testing.
			\item Establishing engagement objectives, scope, and criteria.
			\item Drafting the final audit report draft.
			\item Conducting the formal exit conference.
		\end{enumerate}
	\end{frame}
	
	\begin{frame}[t]{Q51 -- Solution \& Explanation}
		\small
		\textbf{Correct Answer: B} \\[4pt]
		\textbf{Solution:} Standard 2200/2201 states engagement planning begins with objectives, scope, and criteria.
	\end{frame}
	
\end{document}